\reading{CR-23}{An Introduction to Securitisation}{STRIKE FIRST}{7}

\begin{crkeybox}[Learning objectives for this reading]
\begin{enumerate}[label=\textbf{\alph*.},leftmargin=2em]
\item Define securitization, describe the securitization process, and explain the role of participants in the process.
\item Explain the terms over-collateralization, first-loss piece, equity piece, and cash waterfall within the securitization process.
\item Analyze the differences in the mechanics of issuing securitized products using a trust versus a special purpose vehicle (SPV) and distinguish between the three main SPV structures: amortizing, revolving, and master trust.
\item Explain the reasons for and the benefits of undertaking securitization.
\item Describe and assess the various types of credit enhancements.
\item Explain the various performance analysis tools for securitized structures and identify the asset classes they are most applicable to.
\item Define and calculate the delinquency ratio, default ratio, monthly payment rate (MPR), debt service coverage ratio (DSCR), the weighted average coupon (WAC), the weighted average maturity (WAM), and the weighted average life (WAL) for relevant securitized structures.
\item Explain the prepayment forecasting methodologies and calculate the constant prepayment rate (CPR) and the Public Securities Association (PSA) rate.
\end{enumerate}
\end{crkeybox}

\begin{crtrapbox}[Single-layer reading]
CR-23 exists only in the Crash layer.
\end{crtrapbox}

% =====================================================================
\lo{CR-23 a}{Define securitization, describe the securitization process, and explain the role of participants in the process.}

\begin{crdefbox}[Securitization]
Securitization is the process of transforming the \term{illiquid} assets of a
financial institution or corporation into a package of asset-backed securities
(ABSs) or mortgage-backed securities (MBSs).
\end{crdefbox}

\begin{center}
\footnotesize
\begin{tabular}{L{3.6cm} L{10.9cm}}
\toprule
\thead{Participant} & \thead{Role} \\
\midrule
\term{Issuer (SPV/SPE)} & Buys the assets and issues the securities \\
\term{Originator} & Sells the assets to the SPV \\
\term{Investors} & Buy the securities issued by the SPV \\
\term{Structuring agent / sponsor} & Advises on and designs the security,
  including its maturity, credit rating, and enhancements \\
\term{Trustee} & Ensures protection of investors' interests; monitors asset
  conditions \\
\term{Financial guarantor} & Provides financial guarantees against SPV defaults \\
\term{Custodian} & Manages the collection and distribution of payments from the
  asset pool \\
\term{Credit rating agencies} & Assess and rate the securitized products \\
\bottomrule
\end{tabular}
\end{center}
\figcap{Eight participants here against the six in Malz's list at \textbf{CR-22 d},
which is worth noticing: this reading adds the financial guarantor and separates
the structuring agent from the underwriter. The conflicts of interest are treated
at CR-22 d, not here.}

% =====================================================================
\lo{CR-23 b}{Explain the terms over-collateralization, first-loss piece, equity piece, and cash waterfall within the securitization process.}

Securitized notes are divided into classes based on risk. \term{Senior tranches}
have the highest credit quality and receive payments first; \term{junior tranches}
have lower credit quality and receive payments after the senior tranches.

\begin{center}
\footnotesize
\begin{tabular}{L{3.5cm} L{11.0cm}}
\toprule
\thead{Term} & \thead{Meaning} \\
\midrule
\term{Overcollateralization} & Adding more assets to the pool than initially
  needed. Provides a buffer against defaults in the lower-rated tranches \\
\term{First-loss piece} \newline \term{Equity piece} & The lowest credit quality
  tranche, which absorbs initial losses first. It is often held by the
  \emph{originator} to maintain some risk in the deal, which is why it is also
  called the equity piece \\
\term{Cash waterfall} & The priority of payments to the different tranches.
  Senior tranches are paid first, interest and principal. Subordinate tranches are
  paid \emph{only if coverage tests pass}. In case of default, principal is repaid
  starting with the senior tranches \\
\bottomrule
\end{tabular}
\end{center}

\begin{crtrapbox}[Who holds the first-loss piece, and why it matters twice]
The originator typically retains it, which is the retained-risk signal that makes
the structure sellable. It is also the reason, noted at objective d, that
securitisation does \emph{not} deliver 100\% regulatory capital relief, and the
reason an asset-backed note of a given rating can outperform a corporate bond of
the same rating.
\end{crtrapbox}

% =====================================================================
\lo{CR-23 c}{Analyze the differences in the mechanics of issuing securitized products using a trust versus a special purpose vehicle (SPV) and distinguish between the three main SPV structures: amortizing, revolving, and master trust.}

\subsection*{Trust versus SPV}

\begin{center}
\footnotesize
\begin{tabular}{L{3.2cm} L{11.3cm}}
\toprule
\thead{Arrangement} & \thead{Mechanics} \\
\midrule
\term{Trust} & Utilizes a \textbf{two-entity} structure to increase the distance
  between the originator and the assets, enhancing the \emph{legal isolation} of
  the assets \\
\term{SPV} & Typically involves a \textbf{single entity} holding the assets,
  simplifying the structure but potentially \emph{less isolated} than a trust \\
\bottomrule
\end{tabular}
\end{center}

\subsection*{The three SPV structures}

\begin{center}
\footnotesize
\begin{tabular}{L{2.6cm} L{5.4cm} L{6.5cm}}
\toprule
\thead{Structure} & \thead{Used for} & \thead{How it works} \\
\midrule
\term{Amortizing} & Assets with amortization schedules: mortgages, loans &
  Payments are made like coupons, gradually reducing the principal over time.
  Valuation is based on expected maturity and \textbf{weighted average life (WAL)}
  of the assets \\
\term{Revolving} & Assets with frequent paybacks: credit cards, auto loans &
  Principal payments are used to \emph{buy new, similar assets}, keeping the pool
  revolving. Investors are repaid through controlled amortization or lump sum
  payments \\
\term{Master trust} & Frequent issuers under US and UK law; used by MBS and
  credit card ABS originators &
  Allows \textbf{multiple securitizations} to be issued from the same SPV. The
  originator transfers assets to the master trust SPV, and notes are issued out of
  the asset pool based on investor demand \\
\bottomrule
\end{tabular}
\end{center}
\figcap{The asset column is the fastest discriminator. If the underlying pays down
on a schedule it is amortizing; if it pays back quickly and repeatedly it is
revolving; if the issuer comes to market repeatedly it is a master trust. Note
that the master trust is defined by the \emph{issuer's} behaviour, not the
asset's.}

% =====================================================================
\lo{CR-23 d}{Explain the reasons for and the benefits of undertaking securitization.}

\gapbadge

{\footnotesize\color{FRMGrey}The source notes carry the objective heading at the
foot of a page and no body text whatsoever. The following is written from the GARP
curriculum.}

\begin{crgapbox}[What the objective asks for: reasons, and benefits to each side]
The driving force behind securitisation has been the need for banks to realise
value from the assets on their balance sheet. Three factors might lead an
institution to securitise part of its balance sheet:

\begin{itemize}
\item if revenues received from assets remain roughly unchanged but the size of
assets has decreased, there will be an \textbf{increase in the return on equity}
ratio;
\item the level of \textbf{capital required} to support the balance sheet is
reduced, which can lead to cost savings or allow the institution to allocate
capital to more profitable business; and
\item to obtain \textbf{cheaper funding}: the interest payable on asset-backed
securities is frequently considerably below the level payable on the underlying
loans, creating a cash surplus for the originating entity.
\end{itemize}

\textbf{The three reasons, grouped.} A bank securitises for funding, for balance
sheet capital management, or for risk management and credit risk transfer.

\smallskip
\term{Funding.} Supports rapid asset growth; diversifies the funding mix and
reduces the cost of funding; and reduces maturity mismatches. The cost reduction
works because securitisation \emph{de-links} the credit rating of the originating
institution from the credit rating of the issued notes, so most SPV notes are
rated higher than bonds issued directly by the bank.

\smallskip
\term{Balance sheet capital management.} Provides regulatory capital relief,
economic capital relief, and diversified sources of capital. An SPV is not a bank,
so it is not subject to the Basel rules and needs only the capital economically
required by its assets. Note that an originator does \textbf{not} obtain 100\%
regulatory capital relief, because it retains the first-loss piece.

\smallskip
\term{Risk management.} Once assets are securitised the credit risk exposure is
reduced considerably, and removed entirely if the bank does not retain the
first-loss piece. Securitisation can also remove non-performing assets from the
balance sheet, which removes credit risk, removes negative sentiment, and frees
regulatory capital.

\smallskip
\textbf{Benefits to investors.} Investors can diversify sectors of interest,
access different and sometimes superior risk-reward profiles, and access sectors
otherwise closed to them. Securitised notes frequently offer better risk-reward
performance than corporate bonds of the same rating and maturity, which often
occurs precisely because the originator holds the first-loss piece.
\end{crgapbox}

\begin{crtrapbox}[The cautionary case attached to this objective]
Northern Rock used securitisation to back its growing share of the UK residential
mortgage market and developed an \emph{over-reliance} on securitised funding to
the detriment of its retail deposit base, with disastrous consequences in the 2007
crash. The funding benefit and the funding concentration risk are the same
mechanism.
\end{crtrapbox}

% =====================================================================
\lo{CR-23 e}{Describe and assess the various types of credit enhancements.}

\begin{enumerate}
\item \term{Overcollateralization.} Issuing securities with a principal value
\emph{less} than the total value of the underlying assets. Provides additional
collateral to absorb losses from defaults, shielding investors.
\item \term{Pool insurance.} An insurance policy covering losses of principal in
the collateral pool due to defaults. Offered by composite insurance companies to
enhance the credit rating of the ABS.
\item \term{Subordinating note classes (tranching).} Dividing the asset pool into
different classes of notes, where junior notes (Class B) are subordinate to senior
notes (Class A). Junior note payments are \emph{deferred} until senior notes are
redeemed or certain performance criteria are met, protecting senior investors.
\item \term{Margin step-up.} A mechanism that \emph{increases} the coupon rate
after a specified call date. It provides additional yield to investors past the
call date, but allows issuers to refinance if the rates become uncompetitive.
\item \term{Excess spread.} The difference between the cash inflows from the
underlying assets and the cash outflows from interest payments. Any remaining
funds after covering administrative expenses are held in a reserve account for
future losses, or returned to the originator if not needed.
\end{enumerate}

\begin{crkeybox}[Which are internal and which are external]
Overcollateralization, tranching and excess spread are \term{internal}: they are
built out of the deal's own assets and cash flows. Pool insurance is
\term{external}, since it depends on a third-party insurer and therefore
introduces the counterparty risk discussed at CR-14 i. The margin step-up is not
credit enhancement in the loss-absorbing sense at all; it is a yield mechanism
that compensates investors for extension.
\end{crkeybox}

% =====================================================================
\lo{CR-23 f}{Explain the various performance analysis tools for securitized structures and identify the asset classes they are most applicable to.}

\subsection*{Auto loan ABS performance tools}

\begin{itemize}
\item \term{Loss curve.} Shows expected cumulative losses over the life of the
collateral pool, used to compare actual losses to projections. \textbf{Prime loans
have a flatter loss curve; subprime loans have a steeper one.}
\item \term{Absolute prepayment speed (APS).} Measures the expected maturity of
the ABS. Shows the percentage of the total pool balance repaid in a specific
period. Important for determining the value of the implicit call option.
\end{itemize}

\gapbadge

{\footnotesize\color{FRMGrey}The objective asks which asset classes each tool
applies to. The source notes cover the auto loan tools and then say ``Pls check
the following Learning Objective for CC Tools'', without ever supplying the
asset-class mapping. The following is written from the GARP curriculum.}

\begin{crgapbox}[What the objective asks for: the tool-to-asset-class mapping]
\centering\footnotesize
\begin{tabular}{L{4.6cm} L{4.2cm} L{4.6cm}}
\toprule
\thead{Performance measure} & \thead{Calculation} & \thead{Typical asset class} \\
\midrule
Public Securities Association (PSA) &
  $\bigl[\mathrm{CPR}/(0.2 \times \text{months})\bigr]\times 100$ &
  mortgages, home-equity, student loans \\
Constant prepayment rate (CPR) & $1-(1-\mathrm{SMM})^{12}$ &
  mortgages, home-equity, student loans \\
Single monthly mortality (SMM) & prepayment / outstanding pool balance &
  mortgages, home-equity, student loans \\
Weighted average life (WAL) & $\sum (a/365)\cdot \mathrm{PF}(s)$ & mortgages \\
Weighted average maturity (WAM) & weighted maturity of the pool & mortgages \\
Weighted average coupon (WAC) & weighted coupon of the pool & mortgages \\
Debt service coverage ratio (DSCR) & net operating income / debt payments &
  \textbf{commercial} mortgages \\
Monthly payment rate (MPR) & collections / outstanding pool balance &
  \textbf{all non-amortising} asset classes \\
\bottomrule
\end{tabular}

\raggedright
\smallskip
The last two rows are the ones that discriminate. DSCR is a \emph{commercial}
mortgage measure, not a residential one, and MPR applies to all
\emph{non-amortising} classes, which is what makes it the credit card measure.
\end{crgapbox}

% =====================================================================
\lo{CR-23 g}{Define and calculate the delinquency ratio, default ratio, monthly payment rate (MPR), debt service coverage ratio (DSCR), the weighted average coupon (WAC), the weighted average maturity (WAM), and the weighted average life (WAL) for relevant securitized structures.}

\subsection*{Credit card performance tools}

\begin{crfmlbox}[The three credit card ratios]
\[ \text{Delinquency ratio} \;=\;
   \frac{\text{CC receivables more than 90 days overdue}}{\text{Total CC receivables}} \]
\[ \text{Default ratio} \;=\;
   \frac{\text{Written-off CC receivables}}{\text{Total CC receivables}} \]
\[ \mathrm{MPR} \;=\;
   \frac{\text{Monthly principal} + \text{interest payments}}{\text{Total CC receivables}} \]
\vk{All three share the same denominator, total credit card receivables. What
separates them is the numerator: receivables that are \emph{late}, receivables
that are \emph{gone}, and cash that has \emph{arrived}.}
\end{crfmlbox}

\subsection*{MBS performance tools}

\begin{crfmlbox}[Debt service coverage ratio]
\[ \mathrm{DSCR} \;=\; \frac{\text{Net operating income}}{\text{Total debt payments}}
   \;=\; \frac{\text{Income} - \text{Operating expenses}}{\text{Interest} + \text{Principal}} \]
\vk{A DSCR greater than 1 indicates a better ability to pay. Note that the
denominator is interest \emph{and} principal, not interest alone.}
\end{crfmlbox}

\begin{itemize}
\item \term{Weighted average coupon (WAC).} Calculated by multiplying the mortgage
rate for each pool of loans by its loan balance and then dividing by the total
outstanding loan balance for all pools.
\item \term{Weighted average maturity (WAM).} The weighted average months
remaining to maturity for the pool of mortgages in the MBS.
\item \term{Weighted average life (WAL).} Calculated by summing the time to
maturity multiplied by a \term{pool factor}, which is the outstanding notional
value adjusted by the repayment weighting.
\end{itemize}

\begin{crfmlbox}[Weighted average life]
\[ \mathrm{WAL} \;=\; \sum_{s} \left(\frac{a}{365}\right)\cdot \mathrm{PF}(s) \]
\vk{$a$ = actual days in the period; $\mathrm{PF}(s)$ = the pool factor at $s$,
that is, the outstanding notional as a fraction of the original notional. The
$a/365$ term is the actual/365 day count, so each period is weighted by its true
length, not by a nominal quarter.}
\end{crfmlbox}

\begin{crexbox}[Weighted average life of a securitisation]
A note issue with an original notional of \$89,529,500 amortises over nineteen
interest payment dates from 21 November 2003 to 24 July 2008. Principal repayments
run from \$5,058,824 in the first period down to \$3,058,824 in the last,
reducing the outstanding balance to \$16,470,588. Compute the WAL.
\tcblower
Each period contributes $\mathrm{PF}(s) \times (a/365)$, where the pool factor is
the outstanding balance divided by \$89,529,500.

\smallskip
First period: $\mathrm{PF} = 1.00$, $a = 66$ days, contribution
$1.00 \times 66/365 = 0.18082192$.

Second period: $\mathrm{PF} = 84{,}470{,}588/89{,}529{,}500 = 0.943495$,
$a = 91$ days, contribution $0.943495 \times 91/365 = 0.23522739$.

Final period: $\mathrm{PF} = 16{,}470{,}588/89{,}529{,}500 = 0.183965$,
$a = 91$ days, contribution $0.04586606$.

\smallskip
Summing all nineteen contributions:
\[ \mathrm{WAL} \;=\; \mathbf{2.58\ \text{years}} \]

\smallskip
Note what this tells you: the notes have a legal final maturity of roughly 4.7
years, but a weighted average life of only 2.58 years, because principal is
returned throughout the life of the deal. It is this time-weighted maturity that
allows an investor to compare the MBS with other investments of similar maturity,
and the test applies uniquely to MBS precisely because principal comes back early.
\end{crexbox}

% =====================================================================
\lo{CR-23 h}{Explain the prepayment forecasting methodologies and calculate the constant prepayment rate (CPR) and the Public Securities Association (PSA) rate.}

Forecasting prepayments is crucial to computing the cash flows of an MBS. The
underlying payment remains unchanged, but prepayments, for a given price,
\emph{reduce the yield} on the MBS.

\begin{enumerate}
\item \term{Conditional (or constant) prepayment rate (CPR)}: an
\emph{annualized} prepayment rate.
\item \term{Single monthly mortality rate (SMM)}: converts CPR to a
\emph{monthly} prepayment rate.
\end{enumerate}

\begin{crfmlbox}[CPR, SMM and PSA]
\[ \mathrm{CPR} \;=\; 1-(1-\mathrm{SMM})^{12}
   \qquad\qquad
   \mathrm{SMM} \;=\; 1-(1-\mathrm{CPR})^{1/12} \]
\[ \mathrm{PSA} \;=\; \left[\frac{\mathrm{CPR}}{0.2 \times \text{months}}\right]
   \times 100 \]
\vk{The first two are exact inverses of one another; the direction you need
depends on whether you are given an annual or a monthly figure. The 0.2 in the PSA
formula is the 0.2\% per month by which the benchmark CPR ramps up.}
\end{crfmlbox}

\subsection*{The PSA benchmark}

The PSA prepayment benchmark assumes \emph{increasing} monthly prepayment rates as
a pool ages. The standard benchmark is 100\% PSA, and multiples such as 150 PSA or
200 PSA scale it.

\begin{crkeybox}[The 100 PSA ramp]
$\mathrm{CPR} = 0.2\% \times \text{month}$, for months 1 to 29.

\smallskip
$\mathrm{CPR} = 6\%$, from month 30 onward to month 360.
\end{crkeybox}

\begin{center}
\begin{tikzpicture}
\begin{axis}[
  width=12.2cm, height=5.6cm,
  xlabel={Loan age (months)}, ylabel={CPR (\%)},
  xlabel style={font=\sffamily\scriptsize}, ylabel style={font=\sffamily\scriptsize},
  tick label style={font=\sffamily\scriptsize},
  xmin=0, xmax=48, ymin=0, ymax=13.5,
  xtick={0,5,10,20,30,40,48},
  grid=major, grid style={FRMRule!50, line width=0.3pt},
  legend style={font=\scriptsize\sffamily, draw=FRMRule, fill=white,
    at={(0.5,-0.30)}, anchor=north, legend columns=-1, column sep=10pt}]
\addplot[draw=FRMRed, line width=1.1pt, sharp plot]
  coordinates {(0,0) (30,12.0) (48,12.0)};
\addlegendentry{200 PSA}
\addplot[draw=FRMGold, line width=1.1pt, sharp plot]
  coordinates {(0,0) (30,9.0) (48,9.0)};
\addlegendentry{150 PSA}
\addplot[draw=FRMNavy, line width=1.2pt, sharp plot]
  coordinates {(0,0) (30,6.0) (48,6.0)};
\addlegendentry{100 PSA}
\addplot[only marks, mark=*, mark size=2pt, draw=FRMGold, fill=FRMGold]
  coordinates {(5,1.5)};
\addplot[only marks, mark=*, mark size=2pt, draw=FRMNavy, fill=FRMNavy]
  coordinates {(5,1.0)};
\node[font=\sffamily\scriptsize\bfseries, text=FRMNavy, fill=white, inner sep=1.5pt,
  align=left] at (axis cs:22.5,2.1) {month 5: 1.0\% at 100 PSA,\\1.5\% at 150 PSA};
\draw[-{Stealth[length=4pt]}, FRMGrey] (axis cs:16.0,1.8) -- (axis cs:5.5,1.3);
\draw[FRMGrey, dashed, line width=0.7pt] (axis cs:30,0) -- (axis cs:30,12.0);
\node[font=\sffamily\scriptsize\itshape, text=FRMGrey, fill=white, inner sep=1.5pt]
  at (axis cs:40.0,1.2) {ramp ends at month 30};
\draw[-{Stealth[length=4pt]}, FRMGrey] (axis cs:36.0,1.5) -- (axis cs:30.3,2.6);
\end{axis}
\end{tikzpicture}
\end{center}
\figcap{The PSA ramp. Every multiple is the same shape: a linear climb over the
first 30 months and a flat line thereafter. What the multiple scales is the
\emph{height}, not the length of the ramp --- 200 PSA does not reach its plateau
sooner, it reaches a higher one. The ramp is why a seasoned pool and a new pool
prepay at completely different rates under the same PSA assumption.}

\begin{crexbox}[CPR, PSA multiple, and SMM]
A pool is 5 months old and is assumed to prepay at 150 PSA. Calculate the CPR and
the corresponding SMM.
\tcblower
\textbf{Step 1. The 100 PSA benchmark CPR at month 5.}
\[ \mathrm{CPR} = 5 \times 0.2\% = 1\% \]

\textbf{Step 2. Scale to 150 PSA.}
\[ 1.5 \times 0.01 = 0.015 \text{, that is, } 1.5\% \]

\textbf{Step 3. Convert to a monthly rate.}
\[ \mathrm{SMM} = 1-(1-0.015)^{1/12} = \mathbf{0.001259} \]

\smallskip
\textbf{Check it the other way.} The PSA formula recovers the multiple:
$\bigl[0.015/(0.002\times5)\bigr]\times 100 = 150$. And $1-(1-0.001259)^{12}$
returns 0.015, confirming that the CPR and SMM formulas are exact inverses.
\end{crexbox}

\begin{crtrapbox}[Three places to lose a mark]
\textbf{Annual versus monthly.} CPR is annualized and SMM is monthly. The
conversion is a twelfth \emph{root}, not a division by twelve: $0.015/12 =
0.00125$, which is close enough to 0.001259 to look right and is wrong.

\smallskip
\textbf{The ramp cut-off.} The linear ramp applies to months 1 to 29 only. Beyond
month 30 the 100 PSA CPR is a flat 6\%, so month 40 and month 300 give the same
answer.

\smallskip
\textbf{Which direction the multiple applies.} The PSA multiple scales the CPR,
not the month. At 150 PSA in month 5 the CPR is 1.5\%, not the month-7.5 value.
\end{crtrapbox}
